
\documentclass{article}
\usepackage{colortbl}
\usepackage{makecell}
\usepackage{multirow}
\usepackage{supertabular}

\begin{document}

\newcounter{utterance}

\centering \large Interaction Transcript for game `cladder', experiment `full\_v1.5\_default', episode 6304 with Qwen3.5{-}9B.
\vspace{24pt}

{ \footnotesize  \setcounter{utterance}{1}
\setlength{\tabcolsep}{0pt}
\begin{supertabular}{c@{$\;$}|p{.15\linewidth}@{}p{.15\linewidth}p{.15\linewidth}p{.15\linewidth}p{.15\linewidth}p{.15\linewidth}}
    \# & \multicolumn{2}{c}{Player} && \multicolumn{2}{c}{Game Master} \\
    \hline

    \theutterance \stepcounter{utterance}  
    & & & \multicolumn{4}{p{0.6\linewidth}}{
        \cellcolor[rgb]{0.9,0.9,0.9}{
            \makecell[{{p{\linewidth}}}]{
                \texttt{\tiny{[P1$\langle$GM]}}
                \texttt{You are an expert in causal inference. The following question is not a typical commonsense query, but rather a meticulously designed question created by a professor specializing in causal inference, intended to assess the students' mastery of the course content.} \\
\texttt{Imagine a self{-}contained, hypothetical world with only the following conditions, and without any unmentioned factors or causal relationships: Zuph has a direct effect on rixq and xevu. Rixq has a direct effect on xevu. The overall probability of rixq is 32\%. The probability of not rixq and xevu is 35\%. The probability of rixq and xevu is 21\%. Is the chance of xevu smaller when observing rixq?} \\
\texttt{Start your answer with 'Yes' or 'No', followed by additional reasoning or evidence to support your explanation.} \\
            }
        }
    }
     \\ \\

    \theutterance \stepcounter{utterance}  
    & \multicolumn{4}{p{0.6\linewidth}}{
        \cellcolor[rgb]{0.561,0.737,0.561}{
            \makecell[{{p{\linewidth}}}]{
                \texttt{\tiny{[P1$\rangle$GM]}}
                \texttt{No} \\
\\ 
\texttt{**Reasoning:**} \\
\\ 
\texttt{To determine whether the chance of \$xevu\$ is smaller when observing \$rixq\$, we must compare the conditional probability \$P(xevu \textbackslash{}mid rixq)\$ with the marginal probability \$P(xevu)\$.} \\
\\ 
\texttt{1.  **Calculate the marginal probability \$P(xevu)\$:**} \\
\texttt{    We are given the probability of the joint event where \$rixq\$ is false and \$xevu\$ is false: \$P(\textbackslash{}neg rixq \textbackslash{}land \textbackslash{}neg xevu) = 0.35\$.} \\
\texttt{    In a binary system (assuming \$rixq\$ and \$xevu\$ are binary variables taking values True/False or 1/0), the sum of all four joint probabilities must equal 1:} \\
\texttt{    \$\$P(rixq \textbackslash{}land xevu) + P(rixq \textbackslash{}land \textbackslash{}neg xevu) + P(\textbackslash{}neg rixq \textbackslash{}land xevu) + P(\textbackslash{}neg rixq \textbackslash{}land \textbackslash{}neg xevu) = 1\$\$} \\
\\ 
\texttt{    We know:} \\
\texttt{    *   \$P(rixq \textbackslash{}land xevu) = 0.21\$} \\
\texttt{    *   \$P(\textbackslash{}neg rixq \textbackslash{}land \textbackslash{}neg xevu) = 0.35\$} \\
\\ 
\texttt{    The sum of these two is \$0.21 + 0.35 = 0.56\$.} \\
\texttt{    Therefore, the sum of the remaining two probabilities is \$1 {-} 0.56 = 0.44\$.} \\
\texttt{    \$\$P(rixq \textbackslash{}land \textbackslash{}neg xevu) + P(\textbackslash{}neg rixq \textbackslash{}land xevu) = 0.44\$\$} \\
\\ 
\texttt{    We are also given the marginal probability of \$rixq\$: \$P(rixq) = 0.32\$.} \\
\texttt{    Since \$P(rixq) = P(rixq \textbackslash{}land xevu) + P(rixq \textbackslash{}land \textbackslash{}neg xevu)\$, we can solve for the unknown joint term:} \\
\texttt{    \$\$0.32 = 0.21 + P(rixq \textbackslash{}land \textbackslash{}neg xevu)\$\$} \\
\texttt{    \$\$P(rixq \textbackslash{}land \textbackslash{}neg xevu) = 0.32 {-} 0.21 = 0.11\$\$} \\
\\ 
\texttt{    Now substitute this back into the sum equation to find \$P(\textbackslash{}neg rixq \textbackslash{}land xevu)\$:} \\
\texttt{    \$\$0.11 + P(\textbackslash{}neg rixq \textbackslash{}land xevu) = 0.44\$\$} \\
\texttt{    \$\$P(\textbackslash{}neg rixq \textbackslash{}land xevu) = 0.33\$\$} \\
\\ 
\texttt{    Finally, calculate the total marginal probability of \$xevu\$:} \\
\texttt{    \$\$P(xevu) = P(rixq \textbackslash{}land} \\
            }
        }
    }
    & & \\ \\

    \theutterance \stepcounter{utterance}  
    & & & \multicolumn{2}{p{0.3\linewidth}}{
        \cellcolor[rgb]{0.9,0.9,0.9}{
            \makecell[{{p{\linewidth}}}]{
                \texttt{\tiny{[GM$|$GM]}}
                \texttt{No} \\
            }
        }
    }
    & & \\ \\

    \theutterance \stepcounter{utterance}  
    & & & \multicolumn{2}{p{0.3\linewidth}}{
        \cellcolor[rgb]{0.9,0.9,0.9}{
            \makecell[{{p{\linewidth}}}]{
                \texttt{\tiny{[GM$|$GM]}}
                \texttt{no} \\
            }
        }
    }
    & & \\ \\

    \theutterance \stepcounter{utterance}  
    & & & \multicolumn{2}{p{0.3\linewidth}}{
        \cellcolor[rgb]{0.9,0.9,0.9}{
            \makecell[{{p{\linewidth}}}]{
                \texttt{\tiny{[GM$|$GM]}}
                \texttt{game\_result = WIN} \\
            }
        }
    }
    & & \\ \\

\end{supertabular}
}

\end{document}
